\begin{table}[htbp]
\centering\small
\caption{Net account-return dispersion on fixed panel cycles (percent).}
\label{tab:r3_returns}
\setlength{\tabcolsep}{3.5pt}
\begin{tabular}{llrrrrrr}
\hline
Panel & Variant & Accounts & Mean & SD & Median & Range & Descriptive interval \\
\hline
decision & B & 13 & -1.79 & 0.24 & -1.74 & [-2.04, -1.57] & [-4.42, -1.24] \\
decision & A & 13 & -1.57 & 0.39 & -1.59 & [-1.95, -1.17] & [-4.42, -0.70] \\
decision & S & 13 & -1.73 & 0.19 & -1.63 & [-1.95, -1.61] & [-3.18, -0.83] \\
decision & AS & 13 & -1.54 & 0.68 & -1.65 & [-2.16, -0.82] & [-4.42, -0.43] \\
decision & P & 13 & -1.45 & 0.34 & -1.43 & [-1.79, -1.11] & [-4.42, -0.58] \\
upstream & B & 9 & 0.28 & 0.14 & 0.37 & [0.12, 0.37] & [-1.15, 2.26] \\
upstream & AS & 9 & 0.62 & 0.22 & 0.75 & [0.37, 0.75] & [-0.51, 2.26] \\
\hline
\end{tabular}
\par\smallskip\footnotesize Each return is a mean over 13 decision-panel or nine upstream-panel symbol--block accounts. Accounts use the same frozen prices/costs and only panel cycles; not a portfolio or the broad 51-account result. Three executions per row. SD uses ddof=1. Intervals resample whole blocks and paired execution indices together (10,000 draws, seed 20261102). Finite BA draws: decision 9612/10,000; upstream 6986/10,000. Other undefined counts are in JSON; return intervals use all draws. Conditional finite-draw percentiles are descriptive, with only three runs/blocks; no calibrated-coverage claim. Panel support differs and is not pooled.
\end{table}
